\section{Finite formulations of the consequences}
\label{app:finite-consequences}

Some of the reductions in Section~\ref{sec:consequences} use a bipartite
Unique Games instance in which a nearly satisfying labeling satisfies
all constraints incident to most left vertices.  We first obtain a simple
graph with this property and with a prescribed lower bound on its common
left degree.  We then use it to give the finite rational formulation of
the ordering consequence.  The local ordering construction follows
Guruswami, H\aa stad, Manokaran, Raghavendra, and
Charikar~\cite{GuruswamiEtAl2011OrderingCSP} and uses the
low-influence results cited below.  The composition specifies the graph,
the independent neighbor choices, and the conversion to constraints on distinct
variables.

\subsection{A simple strong-left form of Unique Games}

Orient the constraints of a bipartite game from its left side to its
right side.  A left vertex is \emph{strongly satisfied} by a labeling
when every edge incident to it is satisfied.

\begin{lemma}[Simple strong-left games]\label{lem:simple-strong-left}
Fix \(0<\zeta,\nu<1\) and integers \(d,H\ge1\).  There are integers
\(h,R\ge1\), with \(h\ge H\) and \(d\mid h\), and a deterministic
polynomial-time reduction from \(\mathrm{3SAT}\) to nonempty explicit
unweighted simple bipartite Unique Games instances
\(\Phi_\varphi=(A,B,E)\) over \([R]\) with the following properties.
Every left vertex has degree \(h\).  If \(\varphi\) is satisfiable, one
labeling strongly satisfies at least a \(1-\zeta\) fraction of \(A\).
If \(\varphi\) is unsatisfiable, then \(\val(\Phi_\varphi)<\nu\).
The degree \(h\) and the two errors used in
Theorem~\ref{thm:ugc} may be chosen from \(\zeta,\nu,d,H\) before that
theorem supplies the alphabet \(R\).
\end{lemma}

\begin{proof}
We adapt the normalization and sequence construction of Khot and Regev
\cite[Lemmas~3.3, 3.4, and~3.6]{KhotRegev2008}.  The extra copies of
right vertices used below ensure that the final graph is simple.

Choose a multiple \(h\) of \(d\) with
\[
 h\ge\max\{H,2\},\qquad \frac1h<\frac{\nu}{4}.
\]
For \((T)_h=T(T-1)\cdots(T-h+1)\), choose a fixed integer \(T\ge h\)
such that
\[
 p_{\mathrm{graph}}:=\frac{(T)_h}{T^h}>\frac34.
\]
Such a choice exists because \(1-(T)_h/T^h\le\binom h2/T\).
Choose positive rationals \(\beta,\gamma,u\) and a positive integer
\(L_2\) so that
\[
 \beta<\frac{\nu}{4},\qquad
 h^2\gamma<\frac{\nu}{4},\qquad
 (h+1)u<\frac{\zeta}{2},\qquad
 \frac1{L_2}<\min\left\{\frac{\beta\gamma}{2},
                            \frac{\zeta}{2h}\right\}.
\]
Finally choose positive rational \(\gamma_0,\zeta_0<1/2\) with
\[
 2\gamma_0<\frac{\beta\gamma}{2},
 \qquad \sqrt{2\zeta_0}<u.
\]
All these choices precede the invocation of
Theorem~\ref{thm:ugc} with completeness error \(\zeta_0\) and soundness
error \(\gamma_0\).

Write its output as \(\Phi_0=(X_0,Y,E_0)\), orienting each permutation
left-to-right and renumbering the fixed alphabet as \([R]\).
Put \(M_0=|E_0|>0\), give each edge weight \(1/M_0\), and let \(w_x\)
be the total incident weight at \(x\).  Simplicity gives one permutation
at each pair of positive weight.  Create
\[
 n_x=\lfloor 2|X_0|w_x\rfloor
\]
copies of \(x\), each with local edge weights \(w_{xy}/w_x\).
Only vertices with \(w_x>0\) have copies.  If \(X_1\) is the copy set,
then
\[
 |X_0|\le |X_1|=\sum_x n_x\le2|X_0|,
\]
and the incident weight at each copy is one.

Here and below a local satisfied weight uses the labels at that left
vertex and at the shared right vertices.  Given any labeling of the
copy game, keep the right labels and give all copies of each \(x\) a
label that maximizes their common local objective.  The corresponding
labeling \(L\) of \(\Phi_0\) satisfies
\[
 \sum_{x'\in X_1} w_{\mathrm{sat}}(x')
 \le
 \sum_{x:n_x>0}\frac{n_x}{w_x}w_{\mathrm{sat},L}(x)
 \le 2|X_0|\,w_{\mathrm{sat},L}(\Phi_0).
\]
Consequently the copy game's normalized value is at most
\(2\gamma_0\) on a NO input.  On a YES input, lift a labeling of
\(\Phi_0\) with unsatisfied weight at most \(\zeta_0\).  The same
calculation for unsatisfied weight gives mean local loss at most
\(2\zeta_0<u^2\).  Markov's inequality leaves more than a \(1-u\)
fraction of \(X_1\) with local satisfied weight at least \(1-u\).

To remove the weights, put \(\alpha=L_2|Y|\).  At each \(x\in X_1\),
write \(\widetilde w_{xy}\) for its normalized local weights.
Choose one neighbor \(y_0\) of positive local weight, round
\(\alpha\widetilde w_{xy}\) down for every \(y\ne y_0\), and assign the remaining
occurrences to \(y_0\).  All occurrences at a pair retain that pair's
permutation.  This gives an unweighted multigraph
\(\Phi_{\mathrm{occ}}\) of left degree \(\alpha\).  For any subset of
the neighbors, the difference between
its old weight and its new occurrence fraction has magnitude at most
\(|Y|/\alpha=1/L_2\): sum the rounding losses over the subset if it
omits \(y_0\), and over its complement otherwise.  This applies to the
satisfied subset for every labeling.  Hence
\begin{equation}
 \val(\Phi_{\mathrm{occ}})\le 2\gamma_0+\frac1{L_2}<\beta\gamma
 \quad\text{on a NO input}.
\label{eq:strong-left-base-soundness}
\end{equation}
On a YES input, the labeling above has satisfied occurrence fraction
at least \(1-a\), where \(a=u+1/L_2\), at more than a \(1-u\) fraction
of \(X_1\).

Index the \(\alpha\) occurrences at \(x\) by \(j\), with right endpoint
\(y_j\) and permutation \(\pi_j\).  Replace each \(y\in Y\) by
\((y,t)\), \(t\in[T]\), and replace occurrence \(j\) by all \(T\)
occurrences from \(x\) to \((y_j,t)\) with permutation \(\pi_j\).
Call the resulting left-\(\alpha T\)-regular multigraph
\(\Phi_{\mathrm{tag}}\).  For any labeling of \(\Phi_{\mathrm{tag}}\),
keep its left labels and, independently for each original \(y\), give
\(y\) the label of a uniformly chosen copy \((y,t)\).  The expected
edge value in \(\Phi_{\mathrm{occ}}\), counting its occurrences, is
exactly the given edge value in \(\Phi_{\mathrm{tag}}\).  Thus
\begin{equation}
 \val(\Phi_{\mathrm{tag}})\le\val(\Phi_{\mathrm{occ}})<\beta\gamma
 \quad\text{on a NO input}.
\label{eq:strong-left-tag-value}
\end{equation}
The tag-independent lift of the YES labeling preserves every local
satisfied fraction.

Consider the full sequence graph whose left vertices are
\[
 \bigl(x,(j_1,t_1),\ldots,(j_h,t_h)\bigr),
 \qquad j_i\in[\alpha],\quad t_i\in[T],
\]
with position \(i\) joined to \((y_{j_i},t_i)\) by \(\pi_{j_i}\).
This is the sequence construction of Khot and Regev; its left vertex
identifier includes the entire ordered occurrence sequence.  We record
the soundness estimate in a form that permits the left label to depend
on that entire sequence.
Each \(x\) contributes exactly \((\alpha T)^h\) full sequence vertices,
all of degree \(h\).  Thus the uniform edge value is the average satisfied
fraction over a uniform \(x\), an ordered \(h\)-tuple of independent
uniform occurrences at \(x\), and a uniform position in that tuple.

Fix arbitrary right labels in \(\Phi_{\mathrm{tag}}\).  Each occurrence
at \(x\) determines the unique left label that satisfies it.  Let \(p_x(\ell)\)
be the fraction determining label \(\ell\).  The \(p_x(\ell)\) sum to
one, and choosing the best left label at each \(x\) gives
\[
 \mathbb E_x\max_\ell p_x(\ell)\le\val(\Phi_{\mathrm{tag}})<\beta\gamma.
\]
Fewer than a \(\beta\) fraction of \(x\) therefore have
\(\max_\ell p_x(\ell)\ge\gamma\).  At every other \(x\), the probability
that two labels determined by \(h\) independent occurrences agree is
at most
\[
 \binom h2\sum_\ell p_x(\ell)^2
 \le \binom h2\max_\ell p_x(\ell)<h^2\gamma.
\]
If no two determined labels agree, any label of the sequence vertex
satisfies at most one of its \(h\) edges.  On the remaining sequences
its satisfied fraction is at most one.  Averaging first over sequences
and then over \(x\), and using \(1/h+h^2\gamma<1\), bounds the full
sequence graph's value by
\begin{equation}
 s:=\beta+\frac1h+(1-\beta)h^2\gamma<\frac{3\nu}{4}.
\label{eq:strong-left-full-value}
\end{equation}
The estimate holds for every choice of the right labels and every
sequence-dependent choice of left labels.

Retain only sequence vertices for which \(t_1,\ldots,t_h\) are
pairwise distinct.  For each \(x\) and each base sequence
\((j_1,\ldots,j_h)\), exactly \((T)_h\) tag sequences remain.
Each \(x\) thus contributes \(\alpha^h(T)_h\) retained vertices, a
fraction \(p_{\mathrm{graph}}\) of its full set.  The \(h\) right
neighbors of each retained vertex are distinct because their tags are
distinct.  Different occurrence sequences remain different left
vertices, so this retained graph is simple and has common left degree
\(h\).

Give each sequence vertex based at \(x\) the old label of \(x\), and
keep the tag-independent right labels.  For this lifted YES labeling,
satisfaction at position \(i\) depends only on \(j_i\).  Conditioning
the tags to be distinct leaves the base occurrence sequence independent
and uniform.  If the satisfied fraction at \(x\) in
\(\Phi_{\mathrm{occ}}\) is \(f_x\), exactly an \(f_x^h\) fraction of
its retained sequence vertices are strongly satisfied.
Since every \(x\) contributes equally, the strongly satisfied fraction
is at least
\[
 (1-u)(1-a)^h
 \ge 1-u-ha
 =1-(h+1)u-\frac h{L_2}>1-\zeta.
\]
Here \(0<a<1\) follows from the parameter choices.

For NO, extend any labeling of the retained graph to the full graph by
giving discarded left vertices arbitrary labels and keeping all right
labels.  All left degrees are \(h\), so its full edge value is at least
\(p_{\mathrm{graph}}\) times its retained edge value.  By
\eqref{eq:strong-left-full-value}, the latter is at most
\[
 \frac{s}{p_{\mathrm{graph}}}<\frac{3\nu/4}{3/4}=\nu.
\]
This extension permits arbitrary dependence of the retained left
labels on their tags.

All the operations are finite and deterministic.  If \(N\) bounds the
explicit size of \(\Phi_0\), then \(|X_1|=O(N)\) and
\(\alpha=L_2|Y|=O(N)\).  The retained graph has
\(|X_1|\alpha^h(T)_h=O(N^{h+1})\) left vertices and \(h\) edges per
left vertex.  Its identifiers and the rational arithmetic used for
rounding have polynomial bit length, and its permutation tables use
the fixed alphabet inherited from \(\Phi_0\).  This proves the claimed
polynomial bound and completes the reduction.
\end{proof}

\subsection{The finite ordering transfer}

The ordering theorem of Guruswami et al.\ uses a multiscale local
instance and a low-influence rounding estimate.  Their general
outer composition is described as analogous to the two-query
construction in their Section~7.  The following proposition gives
the composition for the rational occurrence-list model of
Section~\ref{sec:consequences}.  It separates the constants chosen
for the ordering gap from the subsequently chosen Unique Games
alphabet.

\begin{proposition}[Finite ordering transfer]
\label{prop:finite-ordering-transfer}
Fix \(k\ge2\) and \(\varnothing\ne\Pi\subsetneq S_k\), put
\(\rho_\Pi=|\Pi|/k!\), and fix
\(0<\delta<(1-\rho_\Pi)/2\).  There are integers
\(q,H\ge1\) and positive rationals \(\zeta,\nu<1\), depending only
on \(k,\Pi,\delta\), with the following property for every fixed
integer \(R\ge1\).

Let \(\Phi=(A,B,E)\) be a nonempty explicit unweighted simple
bipartite Unique Games instance over \([R]\), oriented from \(A\) to
\(B\), whose common left degree \(h\) satisfies \(h\ge H\) and
\(q\mid h\).  There is a deterministic transformation, polynomial in
the explicit size of \(\Phi\) for fixed \(R\), producing an explicit
list of \(\Pi\)-ordering constraints on distinct variables within
each constraint, with nonnegative rational weights, each encoded by
a binary numerator and a positive binary denominator, summing to one.
If some labeling strongly satisfies at least
\((1-\zeta)|A|\) left vertices, then \(\OPT_\Pi(I)\ge1-\delta\).
If \(\val(\Phi)\le\nu\), then
\(\OPT_\Pi(I)\le\rho_\Pi+\delta\).
The constants \(q,H,\zeta,\nu\) are independent of both \(R\) and
\(h\).
\end{proposition}

The output variables will be \((b,z)\in B\times[m]^R\), for a fixed
local symbol set \([m]\) chosen below. At each \(a\in A\), restrict a
global order to \(N(a)\times[m]^R\) and divide that local sorted list
into \(q\) equal consecutive blocks. Averaging block membership over
neighbors, with their edge permutations acting on \(z\), gives one
simplex-valued function of \(z\), whose components have mean \(1/q\)
under uniform \(z\). For fixed \(a\), conditional on the local query
strings, independent neighbor choices make the block ranks independent
with probabilities given by this function. The local estimate bounds
their expected coarse payoff under small influences; the collision
estimate allows correlated strings and bounds the chance of a shared
block, where coarsening can change the induced order. A large influence
in an averaged rank interval instead yields a game label, and many such
vertices contradict low edge value. We then condition on distinct
neighbors to obtain valid constraints and refine the coarse completeness
map to a strict order.

For \(P=\mathbf1_\Pi\), choose \(\pi_0\in\Pi\) and set
\[
 P_0(\sigma)=P(\pi_0\circ\sigma),\qquad \sigma\in S_k.
\]
Then \(P_0(\mathrm{id})=1\), and its average is \(\rho_\Pi\).
For a vector of integer ranks \(r=(r_1,\ldots,r_k)\), let
\(\operatorname{Ref}(r)\) be the permutations compatible with their
weak order: they list positions in increasing rank and may order
positions of equal rank arbitrarily.  Define
\begin{equation}
 \overline{P_0}(r)=
 \frac1{|\operatorname{Ref}(r)|}
 \sum_{\sigma\in\operatorname{Ref}(r)}P_0(\sigma).
\label{eq:ordering-tie-payoff}
\end{equation}
This is the uniform tie extension used in
\cite[Section~8.2]{GuruswamiEtAl2011OrderingCSP}.  It lies in
\([0,1]\) and agrees with \(P_0\) when the ranks are distinct.

The normalization preserves the predicate after one permutation of
tuple positions.  Given a raw tuple \(w=(w_1,\ldots,w_k)\), define
the emitted tuple \(v\) by
\begin{equation}
 v_{\pi_0(j)}=w_j\qquad(j\in[k]).
\label{eq:ordering-slot-permutation}
\end{equation}
If a strict order induces \(\sigma\) on \(w\), it induces
\(\pi_0\circ\sigma\) on \(v\).  More generally, relabeling positions
gives the bijection
\[
 \operatorname{Ref}\bigl(f(v_1),\ldots,f(v_k)\bigr)
 =
 \{\pi_0\circ\sigma:
       \sigma\in\operatorname{Ref}(f(w_1),\ldots,f(w_k))\}
\]
for every rank map \(f\).  Uniform averages over these sets therefore
give equal \(P\)- and \(P_0\)-payoffs, including ties.  The permutation
of positions in \eqref{eq:ordering-slot-permutation} preserves
distinctness and weights.

For a finite local instance \(J=(U,\mathcal C,\omega)\), write
\(t_c=(u_{c,1},\ldots,u_{c,k})\) for the tuple of occurrence
\(c\in\mathcal C\) and \(\omega_c\) for its probability weight.
Its \(q\)-ordering value is
\[
 \operatorname{opt}_q(J)=
 \max_{g:U\to[q]}
 \sum_{c\in\mathcal C}\omega_c\,
 \overline{P_0}\bigl(g(u_{c,1}),\ldots,g(u_{c,k})\bigr).
\]

\begin{lemma}[Rational local ordering data]\label{lem:ordering-local-data}
For every integer \(q\ge1\) and rational \(\eta>0\), there are a
finite local instance \(J=(U,\mathcal C,\omega)\), with
\(|U|=m\ge2\), and probability laws \(\mu_c\) on \([m]^k\) with the
following properties.  Every tuple \(c\) has \(k\) distinct variables;
the weights \(\omega_c\) and all atoms of \(\mu_c\) are rational.
Every one-position marginal of \(\mu_c\) is uniform on \([m]\), and
there is a fixed \(0<\alpha\le1/2\) such that
\(\mu_c(x)>\alpha\) for every \(c,x\).  Moreover
\begin{equation}
 \operatorname{opt}_q(J)\le\rho_\Pi+\eta,\qquad
 v_{\mathrm{loc}}:=
 \sum_c\omega_c\mathbb E_{X\sim\mu_c}\overline{P_0}(X)
 \ge1-\eta.
\label{eq:ordering-local-gap}
\end{equation}
The local laws have common feasible Gram data for the relaxation in
\cite[Section~8.3]{GuruswamiEtAl2011OrderingCSP}: there are a unit
vector \(I\) and vectors \(b_{u,j}\), \(u\in U,j\in[m]\), with
\[
 \sum_j b_{u,j}=I,\quad
 \langle b_{u,i},b_{u,j}\rangle=0\ (i\ne j),\quad
 \|b_{u,j}\|^2=\frac1m,
\]
and, for every two positions \(r,s\) of \(c\),
\begin{equation}
 \langle b_{u_{c,r},i},b_{u_{c,s},j}\rangle
 =\Pr_{X\sim\mu_c}[X_r=i,\ X_s=j].
\label{eq:ordering-local-moments}
\end{equation}
All these data are fixed by \(k,\Pi,q,\eta\).
\end{lemma}

\begin{proof}
The explicit increasing-tuple construction in
\cite[Theorem~11.1, Definition~11.2, and the proof of
Theorem~2.5]{GuruswamiEtAl2011OrderingCSP}, applied to \(P_0\),
gives a finite instance \(J^*\) with
\(\operatorname{opt}_q(J^*)\le\rho_\Pi+\eta\) and a linear order
satisfying every constraint.  Its distribution uses only uniform
choices from finite sets, so its weights are rational.  Its tuple
positions lie in successive disjoint intervals and are therefore
distinct.

Use the finite-copy smoothing from the proof of
\cite[Lemma~8.9]{GuruswamiEtAl2011OrderingCSP}, specialized to this
complete order.  Take \(L\) disjoint copies of \(J^*\), each of weight
\(1/L\), and place their complete orders consecutively in an order
\(\sigma:U\to[m]\), where \(m=L|U^*|\).  The resulting instance \(J\)
still has \(\operatorname{opt}_q(J)\le\rho_\Pi+\eta\), since the value
of any \(q\)-assignment is the average of its values on the copies.
Choose \(L\) and a positive rational \(s_0<1\) so that
\(1/L+s_0<\eta\).  Define a random assignment \(X:U\to[m]\) as
follows: with probability \(1-s_0\), choose a uniform cyclic shift
of the ranks of \(\sigma\); with probability \(s_0\), assign the
variables independently and uniformly from \([m]\).

Each variable is uniform in either part of the mixture.  A cyclic
cut can split at most one of the \(L\) consecutive copies, so every
shift has payoff at least \(1-1/L\).  The mixture's payoff is
therefore at least
\((1-s_0)(1-1/L)>1-\eta\).  Let \(\mu_c\) be its restriction to the
variables of \(c\).  If \(n_c(x)\) is the number of the \(m\) cyclic
shifts with restriction \(x\), then
\[
 \mu_c(x)=(1-s_0)\frac{n_c(x)}m+\frac{s_0}{m^k}.
\]
These atoms are rational and exceed
\(\alpha:=s_0/(2m^k)\).  The denominator \(m^k\) is valid because
the variables in \(c\) are distinct.

Finally regard the finite mixture as a probability space and take
\(I=\mathbf1\) and \(b_{u,j}=\mathbf1_{\{X(u)=j\}}\) in its
\(L^2\) space.  Uniform marginals, disjoint events for different
labels of one variable, and restrictions to each \(c\) give all the
displayed Gram identities.  This constructs the required feasible
data for this particular local instance.
\end{proof}

For scalar functions on the uniform product \([m]^R\), use
\[
 \operatorname{Inf}_\ell(F)
 =\mathbb E\!\left[\operatorname{Var}_{z_\ell}
       (F\mid z_{-\ell})\right].
\]
The noise operator \(T_\theta\), \(0\le\theta\le1\), independently
keeps each coordinate with probability \(\theta\) and otherwise
resamples it uniformly, then averages the function.  In an orthonormal
product expansion these are
\begin{equation}
 T_\theta F=\sum_\sigma\theta^{|\sigma|}\widehat F(\sigma)\chi_\sigma,
 \qquad
 \operatorname{Inf}_\ell(F)
 =\sum_{\sigma:\sigma_\ell\ne0}\widehat F(\sigma)^2.
\label{eq:ordering-noise-influence}
\end{equation}
In particular, further smoothing decreases every influence.

\begin{lemma}[Uniform collision estimate]\label{lem:ordering-collision}
For every \(0<\kappa<1/2\), there is \(\mu_0(\kappa)>0\).  For
each \(m\ge2\), there is a nonnegative function
\(r_c^{m,\kappa}(t)\to0\) as \(t\downarrow0\) with the following
property, uniformly over all \(R\ge1\).  Let
\(F,G:[m]^R\to[0,1]\) have the same mean
\(0<\mu\le\mu_0(\kappa)\) and satisfy
\[
 \operatorname{Inf}_\ell(T_{1-\kappa}F)\le t,\qquad
 \operatorname{Inf}_\ell(T_{1-\kappa}G)\le t
 \quad(\ell\in[R]).
\]
For every coupling \((Z,Z')\) whose two marginals are uniform on
\([m]^R\),
\begin{equation}
 \mathbb E\bigl[
   (T_{1-2\kappa}F)(Z)(T_{1-2\kappa}G)(Z')\bigr]
 \le \mu^{1+\kappa/2}+r_c^{m,\kappa}(t).
\label{eq:ordering-collision}
\end{equation}
The function \(r_c^{m,\kappa}\) is independent of \(R,\mu,F,G\)
and the coupling.
\end{lemma}

\begin{proof}
This is the dimension-uniform form of
\cite[Lemmas~3.4 and~3.5]{GuruswamiEtAl2011OrderingCSP} needed here.
For \(H=T_{1-\kappa}F\), the Fourier formula gives
\[
 \|T_{1-2\kappa}F\|_2^2
 \le \operatorname{Stab}_{1-\kappa}(H),
\]
where \(\operatorname{Stab}_\theta(H)
=\sum_\sigma\theta^{|\sigma|}\widehat H(\sigma)^2\).
Indeed,
\((1-\kappa)^3-(1-2\kappa)^2
=\kappa(1-\kappa-\kappa^2)>0\).
Each full coordinate influence bound for \(H\) implies the corresponding
low-degree influence bound in Theorem~4.4 of Mossel, O'Donnell, and
Oleszkiewicz~\cite{MOO10}.  Apply that theorem with product dimension
\(R\), one-coordinate atom bound \(1/m\), and correlation \(1-\kappa\).
For all sufficiently small \(t\), it gives
\[
 \operatorname{Stab}_{1-\kappa}(H)
 \le \Gamma_{1-\kappa}(\mu)
 +C\frac{\log(2m)}{\kappa}
       \frac{\log\log(1/t)}{\log(1/t)}.
\]
The cutoff and this error are independent of \(R\) and \(\mu\).
Here \(\Gamma_\theta(\mu)\) is the Gaussian halfspace noise stability
of a set of Gaussian measure \(\mu\).

For fixed \(\kappa\), Theorem~B.5 of the same source gives
\(\Gamma_{1-\kappa}(\mu)=\mu^{2/(2-\kappa)+o(1)}\) as
\(\mu\downarrow0\).  Since
\[
 \frac{2}{2-\kappa}-\left(1+\frac{\kappa}{2}\right)
 =\frac{\kappa^2}{2(2-\kappa)}>0,
\]
there is a cutoff \(\mu_0(\kappa)>0\), independent of \(m,R\), below
which \(\Gamma_{1-\kappa}(\mu)\le\mu^{1+\kappa/2}\).
This proves the corresponding squared \(L^2\) bound for both \(F\)
and \(G\), with an error \(r_c^{m,\kappa}(t)\to0\) independent of
the stated variables.  Choose a cutoff \(t_0<e^{-e}\) within the
range of the displayed bound, use its nonnegative error for
\(0<t<t_0\), and set \(r_c^{m,\kappa}(t)=1\) for \(t\ge t_0\).
The latter range follows trivially from the product being at most one.
Cauchy--Schwarz under any coupling with the
specified marginals proves \eqref{eq:ordering-collision}.
\end{proof}

Write
\(\mathcal S_q=\{p\in[0,1]^q:\sum_jp_j=1\}\) and define the
multilinear extension of \(\overline{P_0}\) on \([q]^k\) by
\begin{equation}
 P_0^{\mathrm{ml}}(p_1,\ldots,p_k)
 =\sum_{r\in[q]^k}\overline{P_0}(r)
       \prod_{i=1}^k p_i^{r_i},
 \qquad p_i\in\mathcal S_q.
\label{eq:ordering-multilinear}
\end{equation}

\begin{lemma}[A simplex-valued local estimate]\label{lem:ordering-rounding}
Let \(J,\mu_c,\alpha\) be the data of
Lemma~\ref{lem:ordering-local-data}, and fix \(0<\lambda<1\).
For an integer \(R\ge1\), choose \(c\) with probability \(\omega_c\).
Conditional on this one choice of \(c\), choose the \(R\) coordinate
tuples independently from \(\mu_c\), producing strings
\(Z_1,\ldots,Z_k\in[m]^R\).
Then obtain \(\widetilde Z_i\) by independently resampling every
position and coordinate uniformly with probability \(\lambda\).
For a single function \(F=(F^1,\ldots,F^q):[m]^R\to\mathcal S_q\),
used for every \(c\) and every position, suppose that
\(0<t<1/(kq)\) and
\[
 \operatorname{Inf}_\ell(T_{1-\lambda}F^j)\le t
 \quad(\ell\in[R],\ j\in[q]).
\]
Then
\begin{equation}
 \mathbb E_{c,Z,\widetilde Z}
 P_0^{\mathrm{ml}}\bigl(F(\widetilde Z_1),\ldots,
                         F(\widetilde Z_k)\bigr)
 \le \operatorname{opt}_q(J)+r_s(kqt),
\label{eq:ordering-rounding}
\end{equation}
where \(r_s(u)\to0\) as \(u\downarrow0\), independently of \(R\).
The function \(r_s\) depends only on the fixed local data,
\(k,q,\lambda\).
\end{lemma}

\begin{proof}
This is the simplex-valued estimate of
\cite[Lemma~10.5]{GuruswamiEtAl2011OrderingCSP}, with its sufficient
scalar influence bounds and their dependence made explicit.  Put
\(H=T_{1-\lambda}F\).  Independence of the noise between positions
and multilinearity turn the left side of
\eqref{eq:ordering-rounding} into
\begin{equation}
 \mathbb E_c\mathbb E_{Z\sim\mu_c^{\otimes R}}
 P_0^{\mathrm{ml}}\bigl(H(Z_1),\ldots,H(Z_k)\bigr).
\label{eq:ordering-noise-factorization}
\end{equation}

For each coordinate, the integral ensemble consists of the indicators
\(\mathbf1_{\{X_i=j\}}\) under \(\mu_c\); its atoms have probability
greater than \(\alpha\).  The common Gram data
\eqref{eq:ordering-local-moments} supply one Gaussian ensemble for all
variables of \(J\) with the same first and second moments on every
constraint, as in \cite[Lemma~10.4]{GuruswamiEtAl2011OrderingCSP}.
Explicitly, for a standard Gaussian vector \(g\) in the span of the
Gram vectors, the variables
\[
 G_{u,j}=\langle I,b_{u,j}\rangle+
   \langle b_{u,j}-\langle I,b_{u,j}\rangle I,g\rangle
\]
have those moments.  Use independent copies of this ensemble in the
\(R\) coordinates.  Each \(H^j\) is evaluated on these ensembles
through its multilinear product expansion; equivalently, express an
orthonormal basis on \([m]\) as linear combinations of the symbol
indicators and use the same polynomial for the integral and Gaussian
ensembles.

For a fixed \(c\), let \(\mathcal T_\theta^c\) retain each whole
coordinate tuple with probability \(\theta\) and otherwise resample it
from \(\mu_c\).  Uniform one-position marginals give, componentwise,
\[
 \mathcal T_{1-\lambda}^c[F^j(Z_i)]
 =(T_{1-\lambda}F^j)(Z_i)=H^j(Z_i).
\]
Thus these are the smoothed components required by the invariance
theorem on the product of the tuple ensembles.  The same marginals
give the exact influence identity below; the superscript indicates the
product measure used to compute influence:
\begin{equation}
 \operatorname{Inf}^{\mu_c^{\otimes R}}_\ell
       \bigl(H^j(Z_i)\bigr)
 =
 \operatorname{Inf}^{[m]^R}_\ell(H^j).
\label{eq:ordering-vector-influence}
\end{equation}
Indeed, conditioning on every coordinate tuple except the \(\ell\)th
leaves the \(\ell\)th symbol at position \(i\) uniform, and the other
symbols at that position have the uniform product law.  For a vector,
influence is the sum of its component influences
\cite[Definition~2.6]{IsakssonMossel2009}.  Thus the vector
\((H^j(Z_i))_{i\in[k],j\in[q]}\) has each coordinate influence at
most \(kqt\); each component has variance at most one.

For the remaining tail hypothesis, let \((F^j(Z_i))_S\),
\(S\subseteq[R]\), be the Hoeffding components on the tuple product
\(\mu_c^{\otimes R}\), and let the superscript \(>d\) retain the terms
with \(|S|>d\). The tuple-noise identity above multiplies component
\(S\) by \((1-\lambda)^{|S|}\). Since the unsmoothed \(F^j\) takes
values in \([0,1]\), for every real \(d>0\),
\[
\begin{aligned}
 \operatorname{Var}\bigl((H^j(Z_i))^{>d}\bigr)
 &=\sum_{|S|>d}(1-\lambda)^{2|S|}
       \bigl\|(F^j(Z_i))_S\bigr\|_2^2\\
 &\le (1-\lambda)^{2d}\operatorname{Var}(F^j(Z_i))
 \le (1-\lambda)^{2d}.
\end{aligned}
\]
All norms and variances here use \(\mu_c^{\otimes R}\). This supplies
the tail bound in \cite[Theorem~3.6]{IsakssonMossel2009}.
These are sufficient hypotheses for the vector-valued
invariance estimate stated in
\cite[Theorem~10.1]{GuruswamiEtAl2011OrderingCSP}.

Let \(\mathsf p:\mathbb R^q\to\mathcal S_q\) be Euclidean projection.
It is nonexpansive.  The function
\[
 \Psi(x_1,\ldots,x_k)=
 P_0^{\mathrm{ml}}(\mathsf p(x_1),\ldots,\mathsf p(x_k))
\]
is Lipschitz with constant at most \(\sqrt{kq}\): on the simplex,
telescoping the positions bounds the payoff difference by
\(\sum_i\|x_i-y_i\|_1\), and projection and Cauchy--Schwarz give the
claim.  Apply the cited invariance estimate to this function, separately
for each \(c\), with vector dimension \(kq\) and influence threshold
\(kqt\).  Its displayed error is at most
\[
 r_s(kqt),\qquad
 r_s(u)=C_{kq}\sqrt{kq}\,
             u^{\lambda/(18\log(1/\alpha))}.
\]
In particular the error is independent of \(R\).

On the integral ensemble, \(H(Z_i)\in\mathcal S_q\), so \(\Psi\)
equals the payoff in \eqref{eq:ordering-noise-factorization}.
On the common Gaussian ensemble, project \(H(G_u)\) to
\(\mathcal S_q\) for every \(u\in U\), and independently draw one
\(q\)-label for each \(u\) with those probabilities.  This is one
random assignment of \(J\).  The variables in every local tuple are
distinct, so its conditional expected payoff at that tuple is the
multilinear payoff \(\Psi\).  Every realized assignment has value at
most \(\operatorname{opt}_q(J)\).  Averaging the invariance comparisons
over \(c\) proves \eqref{eq:ordering-rounding}.
\end{proof}

\begin{proof}[Proof of Proposition~\ref{prop:finite-ordering-transfer}]
We first derive estimates that hold for every \(R\), then choose their
numerical parameters independently of \(R\).  Fix for now
\(0<\kappa<1/2\), \(q\) with \(1/q\le\mu_0(\kappa)\), and local data
from Lemma~\ref{lem:ordering-local-data}.  Use the local experiment of
Lemma~\ref{lem:ordering-rounding} with the actual resampling probability
\(\lambda=2\kappa\).  Also fix auxiliary constants
\(0<t<1/(kq)\) and \(\chi>0\).

For an input graph \(\Phi\), let \(\pi_{a\to b}\) denote the permutation
on edge \(ab\).  Its action on strings is fixed by
\[
 (\pi_{a\to b}z)_{\pi_{a\to b}(\ell)}=z_\ell
 \qquad(\ell\in[R]).
\]
The output variable set will be
\(\mathcal V=B\times[m]^R\).  Consider the following independent
neighbor experiment: choose \(a\) uniformly from \(A\), choose
\(b_1,\ldots,b_k\) independently and uniformly from \(N(a)\), and
independently perform the local experiment, obtaining
\(\widetilde Z_1,\ldots,\widetilde Z_k\).  Its raw output positions are
\begin{equation}
 w_i=(b_i,\pi_{a\to b_i}\widetilde Z_i)\in\mathcal V.
\label{eq:ordering-iid-queries}
\end{equation}
For any rank map \(f:\mathcal V\to\mathbb Z\), define
\[
 V(f)=\mathbb E\,
   \overline{P_0}\bigl(f(w_1),\ldots,f(w_k)\bigr),
 \qquad V_a(f)=\mathbb E[
   \overline{P_0}(f(w_1),\ldots,f(w_k))\mid a].
\]
The bounded payoff \eqref{eq:ordering-tie-payoff} is used even if the
experiment repeats an output ID.  We will obtain a lower bound for one
coarse rank map and an upper bound for every injective order, then
transfer both bounds to distinct-variable constraints.

Suppose a labeling \(L\) strongly satisfies at least \(1-\zeta\) of
the left vertices.  Define the coarse rank map
\(f(b,z)=z_{L(b)}\).  At every strongly satisfied \(a\) and every
neighbor \(b\),
\[
 f(b,\pi_{a\to b}z)=z_{\pi_{a\to b}^{-1}(L(b))}=z_{L(a)}.
\]
Thus all \(k\) positions read the same label coordinate, for every
choice of neighbors.  With probability \((1-2\kappa)^k\) none of
these \(k\) symbols is resampled; conditional on that event their
tuple has law \(\mu_c\).  Nonnegativity of the payoff gives
\begin{equation}
 V(f)\ge (1-\zeta)(1-2\kappa)^k v_{\mathrm{loc}}.
\label{eq:ordering-iid-completeness}
\end{equation}

For soundness fix an injective order \(O:\mathcal V\to[|\mathcal V|]\).
At a fixed \(a\), its local set \(N(a)\times[m]^R\) has exactly
\(M=hm^R\) distinct IDs.  Sort this set by \(O\) and divide the sorted
list into \(q\) consecutive blocks \(C_{a,1},\ldots,C_{a,q}\) of
size \(M/q\).  These sizes are integers because \(q\mid h\).
Each block is the intersection of the local set with an interval of
global ranks.  Define
\begin{equation}
 F_a^j(z)=\frac1h\sum_{b\in N(a)}
  \mathbf1_{\{(b,\pi_{a\to b}z)\in C_{a,j}\}},
 \qquad F_a=(F_a^1,\ldots,F_a^q).
\label{eq:ordering-block-functions}
\end{equation}
For every \(z\), \(F_a(z)\in\mathcal S_q\).  Each edge permutation
is a bijection on strings, so
\begin{equation}
 \mathbb E_z F_a^j(z)=\frac{|C_{a,j}|}{hm^R}=\frac1q.
\label{eq:ordering-block-mass}
\end{equation}

For any interval \(I\) of integer ranks put
\[
 F_a^I(z)=\frac1h\sum_{b\in N(a)}
   \mathbf1_{\{O(b,\pi_{a\to b}z)\in I\}}.
\]
Call \(a\) \emph{pseudorandom at threshold \(t\)} if
\(\operatorname{Inf}_\ell(T_{1-\kappa}F_a^I)\le t\) for every
\(I\) and every \(\ell\in[R]\).  At such an \(a\), this bound holds
for each block function \(F_a^j\).

Let \(j_a(w)\) be the block index of a local ID \(w\).  Conditional
on the local query strings, independence of the neighbor choices
gives the exact identity
\begin{equation}
 \mathbb E_{b_1,\ldots,b_k}
 \overline{P_0}\bigl(j_a(w_1),\ldots,j_a(w_k)\bigr)
 =
 P_0^{\mathrm{ml}}\bigl(F_a(\widetilde Z_1),\ldots,
                         F_a(\widetilde Z_k)\bigr).
\label{eq:ordering-neighbor-factorization}
\end{equation}
The identity includes outcomes with equal neighbors.  Further
smoothing decreases the block influences, so
\(\operatorname{Inf}_\ell(T_{1-2\kappa}F_a^j)\le t\).
Lemma~\ref{lem:ordering-rounding}, applied to the single function
\(F_a\), bounds the expected coarse payoff in
\eqref{eq:ordering-neighbor-factorization} by
\(\operatorname{opt}_q(J)+r_s(kqt)\).

When the \(k\) coarse ranks are distinct, their induced permutation
is the one induced by the fine ranks \(O(w_i)\).  Otherwise the two
payoffs differ by at most one.  This also covers a repeated output ID,
which necessarily produces a coarse tie.  For two positions \(i\ne i'\),
independence of their noise and neighbor choices gives their coarse
tie probability as
\[
 \sum_{j=1}^q
 \mathbb E_c\mathbb E_{Z\sim\mu_c^{\otimes R}}
 \bigl[
  (T_{1-2\kappa}F_a^j)(Z_i)
  (T_{1-2\kappa}F_a^j)(Z_{i'})
 \bigr].
\]
Conditional on \(c\), both strings have uniform product marginals.
Their correlation is permitted by
Lemma~\ref{lem:ordering-collision}.  Use that lemma with the fixed
mean \(1/q\) from \eqref{eq:ordering-block-mass}, and write
\(r_c=r_c^{m,\kappa}\).  The last display is at most
\(q^{-\kappa/2}+q r_c(t)\).  A union bound over pairs therefore gives,
at every pseudorandom \(a\),
\begin{equation}
 V_a(O)\le
 S_{\mathrm{loc}}:=\operatorname{opt}_q(J)
 +\binom{k}{2}\bigl(q^{-\kappa/2}+q r_c(t)\bigr)
 +r_s(kqt).
\label{eq:ordering-pseudorandom-bound}
\end{equation}

We next turn an excess over this bound into a labeling of \(\Phi\).
For each \(b\in B\), restrict \(O\) to
\(\{b\}\times[m]^R\) and compress its ranks to obtain a strict order
\(o_b\) of \([m]^R\).  Let
\[
 \mathcal L_b=\left\{\ell\in[R]:
   \operatorname{Inf}_\ell
    \bigl(T_{1-\kappa}\mathbf1_{\{o_b(z)\in I\}}\bigr)>t/2
   \text{ for some rank interval }I\right\}.
\]
The few-influential-coordinates lemma
\cite[Definition~4.2 and Lemma~4.3]{GuruswamiEtAl2011OrderingCSP}
gives
\begin{equation}
 |\mathcal L_b|\le\frac{400}{\kappa(t/2)^3}
 =\frac{3200}{\kappa t^3}.
\label{eq:ordering-right-candidates}
\end{equation}
This bound is uniform in \(R\).

If \(a\) is not pseudorandom, choose one witnessing interval \(I_a\)
and coordinate \(\ell_a\) with
\(\operatorname{Inf}_{\ell_a}(T_{1-\kappa}F_a^{I_a})>t\).
Noise commutes with coordinate permutations and averaging, and
influence is a convex squared norm.  Hence
\begin{equation}
 \operatorname{Inf}_{\ell_a}(T_{1-\kappa}F_a^{I_a})
 \le\frac1h\sum_{b\in N(a)}
 \operatorname{Inf}_{\pi_{a\to b}(\ell_a)}
 \left(T_{1-\kappa}
       \mathbf1_{\{O(b,z)\in I_a\}}\right).
\label{eq:ordering-influence-average}
\end{equation}
Each term is at most one.  An average greater than \(t\) forces more
than a \(t/2\) fraction of the terms to exceed \(t/2\).
The restriction of a global-rank interval to one right code is an
interval of its relative order, so these neighbors satisfy
\(\pi_{a\to b}(\ell_a)\in\mathcal L_b\).

Give each such \(a\) label \(\ell_a\), and label every other left
vertex arbitrarily.  Independently at each right vertex with
\(\mathcal L_b\ne\varnothing\), choose one uniform label from
\(\mathcal L_b\); use any label when the set is empty.  By
\eqref{eq:ordering-right-candidates}, the expected satisfied fraction
among the edges at each nonpseudorandom \(a\) is greater than
\[
 \frac t2\,\frac{\kappa t^3}{3200}
 =\frac{\kappa t^4}{6400}.
\]
One random right labeling is used for all left vertices.  Linearity
of expectation and the common left degree identify the resulting
average with the ordinary uniform edge fraction.

If \(V(O)>S_{\mathrm{loc}}+\chi\), more than a \(\chi\) fraction of
left vertices are nonpseudorandom: if that fraction is \(g\), the bound
\(0\le V_a(O)\le1\) and \eqref{eq:ordering-pseudorandom-bound} give
\(V(O)\le(1-g)S_{\mathrm{loc}}+g\le S_{\mathrm{loc}}+g\).  The
preceding decoder would then give a labeling of edge value greater than
\(\chi\kappa t^4/6400\).
Consequently
\begin{equation}
 \val(\Phi)\le\nu<\frac{\chi\kappa t^4}{6400}
 \quad\Longrightarrow\quad
 V(O)\le S_{\mathrm{loc}}+\chi\quad\text{for every injective }O.
\label{eq:ordering-iid-soundness}
\end{equation}

It remains to turn the independent experiment into strict ordering
constraints.  Let \(D\) be the event that \(b_1,\ldots,b_k\) are
distinct.  For \(h\ge k\),
\begin{equation}
 p_D:=\Pr(D)=\frac{(h)_k}{h^k},\qquad
 \delta_h:=1-p_D\le\frac{\binom{k}{2}}h.
\label{eq:ordering-distinct-probability}
\end{equation}
The probability is the same at every \(a\), and \(D\) is independent
of the local experiment.  Conditional on \(D\), \(a\) remains uniform
and its ordered distinct neighbor tuple is uniform among its
\((h)_k\) possibilities.  All output IDs are then distinct because
their first coordinates \(b_i\) differ.  Let \(I_0\) denote this finite
weighted distribution of strict \(P_0\) constraints, and write
\(\OPT_{P_0}(I_0)\) for its maximum strict payoff.

Let \(V_D(f)\) denote the expected tie-extended payoff after this
conditioning, and use \(V_{\neg D}(f)\) for conditioning on its
complement.  For any rank map \(f\), coarse or injective,
\begin{equation}
 V(f)=p_DV_D(f)+(1-p_D)V_{\neg D}(f),
 \qquad |V_D(f)-V(f)|\le\delta_h,
\label{eq:ordering-conditioning}
\end{equation}
because all payoffs lie in \([0,1]\).  Apply this first to the coarse
map in \eqref{eq:ordering-iid-completeness}.  Within each finite fiber
of that map, choose one uniform random permutation of the distinct
IDs, independently between fibers, and concatenate the fibers in
rank order.  On every conditioned tuple, restriction of these
permutations to its IDs is uniform within each tied rank class.
Thus the expected strict \(P_0\)-payoff of this one random global
order equals \(V_D(f)\).  Some global order attains at least that
value.  For soundness, an injective \(O\) has distinct ranks on every
conditioned tuple, so its tie-extended payoff there is its strict
\(P_0\)-payoff.  Equations \eqref{eq:ordering-iid-completeness},
\eqref{eq:ordering-iid-soundness}, and \eqref{eq:ordering-conditioning}
therefore give the strict bounds
\begin{equation}
\begin{aligned}
 \mathrm{YES}:\quad&
 \OPT_{P_0}(I_0)\ge(1-\zeta)(1-2\kappa)^k v_{\mathrm{loc}}-\delta_h,\\
 \mathrm{NO}:\quad&
 \OPT_{P_0}(I_0)\le S_{\mathrm{loc}}+\chi+\delta_h
 \quad\text{if }\nu<\chi\kappa t^4/6400.
\end{aligned}
\label{eq:ordering-conditioned-gap}
\end{equation}
These bounds are for the conditioned \(P_0\) tuples.  Applying
\eqref{eq:ordering-slot-permutation} makes every occurrence use the
original predicate \(\Pi\) with exactly the same values.  Call this
emitted instance \(I\).

We now make the choices for the prescribed \(\delta\).  Choose a
positive rational \(\eta\) with \(10\eta<\delta\), and a rational
\(0<\kappa<1/2\) with \(2\kappa k<\eta\).  Put \(K=\binom{k}{2}\).
Choose an integer \(q\) so large that
\[
 \frac1q\le\mu_0(\kappa),\qquad Kq^{-\kappa/2}<\eta.
\]
Choose the data \(J,m,\alpha\) from
Lemma~\ref{lem:ordering-local-data} at error \(\eta\).  The two
analytic error functions now have fixed parameters and are uniform
in \(R\).  Choose a positive rational \(t<1/(kq)\) so small that
\[
 Kq r_c(t)<\eta,\qquad r_s(kqt)<\eta.
\]
Set \(\chi=\eta\).  Choose positive rationals \(\zeta,\nu<1\) with
\[
 \zeta<\eta,\qquad
 \nu<\frac{\eta\kappa t^4}{6400},
\]
and choose an integer \(H\ge\max\{k,q\}\) with \(K/H<\eta\).
These choices depend only on \(k,\Pi,\delta\).  In particular, they
precede both the choice of \(h\) in
Lemma~\ref{lem:simple-strong-left} and the alphabet supplied by
Theorem~\ref{thm:ugc}.

For every \(h\ge H\), \(\delta_h<\eta\).  The local gap
\eqref{eq:ordering-local-gap}, Bernoulli's inequality, and
\eqref{eq:ordering-conditioned-gap} give
\[
 \mathrm{YES}:\quad
 \OPT_\Pi(I)\ge1-\zeta-2\kappa k-\eta-\delta_h
 >1-4\eta>1-\delta,
\]
whereas \eqref{eq:ordering-pseudorandom-bound} gives
\[
 \mathrm{NO}:\quad
 \OPT_\Pi(I)\le S_{\mathrm{loc}}+\chi+\delta_h
 <\rho_\Pi+6\eta<\rho_\Pi+\delta.
\]
This proves the numerical claims simultaneously for every \(R\).

Finally we specify the encoded transformation.  Once \(R\) is fixed,
the choice of \(c\), the product \(\mu_c^{\otimes R}\), and the
independent rational \(2\kappa\)-noise form one fixed finite rational
law \(\theta_R\) on local outcomes \(\xi\).  These are the actual
local outcomes above; the Gaussian variables used for the analytic
comparison are not part of the output.  Renumber \(B\) by
\([|B|]\) and explicitly list \(\mathcal V=B\times[m]^R\).
For each \(a\in A\), each ordered
distinct neighbor tuple, and each \(\xi\in\operatorname{supp}\theta_R\),
list the tuple in \eqref{eq:ordering-iid-queries}, permuted by
\eqref{eq:ordering-slot-permutation}, with weight
\begin{equation}
 \frac{\theta_R(\xi)}{|A|(h)_k}.
\label{eq:ordering-rational-weight}
\end{equation}
The weights sum to one and are nonnegative rationals.  Every listed
tuple has distinct IDs; the same tuple may occur in several entries.
There are \(|B|m^R\) variables and at most
\(|A|(h)_k|\operatorname{supp}\theta_R|\) entries.  For fixed
\(k,m,R,\theta_R\), these counts are polynomial in the explicit graph
size, even when \(h\) varies.  A fixed common denominator for
\(\theta_R\) shows that each weight has
\(O(\log|A|+k\log h)\) bits plus a fixed constant.  The IDs have
\(O(\log|B|)\) bits plus a fixed constant.  Exact enumeration of
these entries is therefore deterministic and polynomial for fixed
\(R\), as asserted.
\end{proof}

The same occurrence-list encoding also justifies the deterministic
random-order guarantee in Section~\ref{sec:consequences}.  Build an
order by exposing one next variable at a time.  If exactly \(s\)
variables of a constraint have been exposed, its conditional success
probability is the number of permutations in \(\Pi\) whose first
\(s\) positions are the roles of those variables in the exposed order, divided by
\((k-s)!\).  It can be computed by inspecting the fixed set \(S_k\).
For each choice of the next variable, sum these probabilities against
the listed rational weights and choose a next variable whose conditional
expectation is at least the average over the choices.
This preserves the initial expectation \(|\Pi|/k!\).  All comparisons
are exact and polynomial in the input bit length: the product of the
weight denominators has bit length at most their total encoded length,
and the factorial denominators have constant size for fixed \(k\).
The variable list is explicit, so the number of prefix steps and
candidate choices is polynomial as well.
